% Checkpoint36 proposed source. All checkpoint35 sources remain unchanged.
\statement{A.120}{Proved}{Every odd coherent class becomes trivial after the positive-factor centralizer constraint is forgotten}
Retain the actual full-GT negative involution \(\alpha\), its original
parameter \(\phi\), canonical word \(u\), and unit \(d\) of A.104--A.118.
Let \(R=N_1\cap O\) be the actual residual of A.116, and put
\[
 E_R=F/R,\qquad P=F^{(2)},\qquad
 G_p=\operatorname{SL}_2(\mathbb Z_p).
\]
Thus A.112 and A.116 give the actual compact derived group
\[
 E_R'=P'\times G_{\rm odd}',\qquad
 G_{\rm odd}'=G_3'\times\prod_{p\geq5}G_p.
 \tag{A.120.1}
\]
For \(e=d\bmod3\in\{1,-1\}\), put
\(\epsilon=(1-e)/2\) and \(R_e=y^\epsilon r\).
Use the ORIGINAL LM7\_PUSH~E.61/E.65 pure frame \(A_r\) and the
specified negative frame from A.111:
\[
 J_e^\alpha=
 \begin{cases}
 s\phi A_r,&e=1,\\
 A_rs,&e=-1,
 \end{cases}
 \qquad
 \kappa_e=\operatorname{Ad}\bigl((J_e^\alpha)^{-1}\bigr)\alpha .
 \tag{A.120.2}
\]
This is an involution on \(E_R'\). Its restriction preserves
the actual positive-factor centralizer
\[
 \begin{gathered}
 T_{R,e}=C_{E_R}(R_e)=T_{2,e}\times T_{{\rm odd},e},\\
 T_{2,e}=(P')^{\rho_e},\qquad
 T_{{\rm odd},e}=C_{G_3'}(J_e)\times
                   \prod_{p\geq5}C_{G_p}(J_e),\\
 J_e=\operatorname{diag}(1,e)J\operatorname{diag}(1,e)^{-1},
 \qquad
 J=\begin{pmatrix}0&-1\\1&-1\end{pmatrix}.
 \end{gathered}
 \tag{A.120.3}
\]
All products and actions are the actual ones proved in A.112 and A.117.
Then
\[
 H^1(C_2,G_{\rm odd}';\kappa_e)=\{*\},
 \qquad
 H^1(C_2,T_{{\rm odd},e};\kappa_e)
       \simeq\prod_{p\ {\rm odd}}C_2 .
 \tag{A.120.4}
\]
The map induced by the actual inclusion of coefficient groups
\[
 H^1(C_2,T_{R,e};\kappa_e)
       \longrightarrow H^1(C_2,E_R';\kappa_e)
\]
is, under the natural product homeomorphisms, exactly
\[
 \begin{gathered}
 H^1(C_2,T_{2,e};\kappa_e)\times
       \prod_{p\ {\rm odd}}C_2
       \longrightarrow H^1(C_2,P';\kappa_e),\\
 ([z_2],\varepsilon_{\rm odd})\longmapsto
             \operatorname{incl}_*([z_2]).
 \end{gathered}
 \tag{A.120.5}
\]
These are continuous maps of pointed compact Hausdorff spaces.
In particular every class with basepoint dyadic coordinate becomes
trivial after this inclusion, regardless of all its odd coordinates.
No coordinate of the specified original coherent class is evaluated.

\begin{proof}
\emph{A.120.1. The exact matrix action of the original negative frame.}
Use the actual signed congruence representation of A.103--A.104 and
\[
 X=I+2E_{12},\qquad Y=I-2E_{21},\qquad
 S=\begin{pmatrix}0&-1\\1&0\end{pmatrix}.
\]
Conjugation by \(S\) interchanges \(X\) and \(Y\). Thus it implements
the original \(s\)-action on each odd factor. Since \(S^2=-I\),
we use this conjugation action; we do not send the order-two element
\(s\) to \(S\) by an asserted homomorphism into \(\operatorname{SL}_2\).

A.104 gives the actual action of \(\alpha\) on each odd factor as
\(\operatorname{Ad}(K_{\phi,p})\), where
\[
       \det K_{\phi,p}=-1,\qquad K_{\phi,p}^2=I.
\]
Write \(A_p=q_p(A_r)\), \(\Phi_p=q_p(\phi)\), and set
\[
 Q_{p,e}=
 \begin{cases}
 S\Phi_pA_p,&e=1,\\
 A_pS,&e=-1,
 \end{cases}
 \qquad
 L_{p,e}=Q_{p,e}^{-1}K_{\phi,p}.
 \tag{A.120.6}
\]
These are the literal words (A.120.2), in their printed order.
Consequently \(\kappa_e\) acts on \(G_p\) by
\(\operatorname{Ad}(L_{p,e})\), with \(\det L_{p,e}=-1\).
No quotient frame has been substituted for the original \(A_r\).

The exact reversing identity proved in A.112 is
\[
                  L_{p,e}J_eL_{p,e}^{-1}=J_e^{-1}.
 \tag{A.120.7}
\]
Its scalar qualification matters. The original frame relation first
identifies the corresponding conjugation actions on the derived
matrix factor. The centralizer of that factor in the full matrix
algebra is scalar: at three it suffices to commute with \(X^3\)
and \(Y^3\), using the integral-domain entry equations.
Thus the two cyclic matrices initially differ by a scalar.
Both have trace \(-1\), which forces that scalar to be one.
This recovers the literal matrix identity (A.120.7).

Conjugate \(L_{p,e}\) by \(\operatorname{diag}(1,e)\).
Writing the resulting matrix as
\(\left(\begin{smallmatrix}a&b\\c&d_0\end{smallmatrix}\right)\),
the equation \(LJ=J^{-1}L\) gives
\[
                         c=a+b,\qquad d_0=-a.
\]
Hence \(\operatorname{tr}L_{p,e}=0\).
Cayley--Hamilton and \(\det L_{p,e}=-1\) give
\[
                            L_{p,e}^2=I.
 \tag{A.120.8}
\]
There is no division by three. In particular the argument applies
at \(p=3\), even though the order-three centralizer algebra is ramified
there; the reversing involution has two distinct eigenvalues modulo
every odd prime.

\emph{A.120.2. A complete integral matrix-cocycle calculation.}
We prove a slightly more general elementary fact.
Let \(\mathcal O=\mathbb Z_p\), \(p\) odd, and let \(L\) be an invertible
trace-zero matrix with \(L^2=\lambda I\), where
\(\lambda\in\mathcal O^\times\). Put \(\kappa=\operatorname{Ad}(L)\).
Then
\[
                 H^1(C_2,\operatorname{SL}_2(\mathcal O);\kappa)
                         =\{*\}.
 \tag{A.120.9}
\]
We use the convention \(z\kappa(z)=1\), with equivalence
\(z\sim g^{-1}z\kappa(g)\).
For a cocycle \(z\), set \(M=zL\). Direct multiplication gives
\[
                M^2=\lambda I,\qquad
                \det M=\det L=-\lambda .
\]
The Cayley--Hamilton equation for \(M\) therefore reduces to
\((\operatorname{tr}M)M=0\). Since \(M\) is invertible,
\(\operatorname{tr}M=0\). Both \(M\) and \(L\) have characteristic
polynomial \(t^2-\lambda\).

We verify integral conjugacy and the determinant correction.
If \(\lambda\) is a square modulo \(p\), successive lifting gives
a unit square root \(a\in\mathcal O^\times\).
The idempotents
\[
                   \frac{I+M/a}{2},\qquad
                   \frac{I-M/a}{2}
\]
split \(\mathcal O^2\) into two rank-one direct summands.
Each has rank one after reduction modulo \(p\), and each is free
because \(\mathcal O\) is local. The same holds for \(L\).
Bases of these summands conjugate both matrices to
\(\operatorname{diag}(a,-a)\), so
\(M=g_0^{-1}Lg_0\) for some \(g_0\in\operatorname{GL}_2(\mathcal O)\).
In this basis the centralizer of \(L\) contains
\(\operatorname{diag}(t,1)\) for every unit \(t\), and therefore
has every possible unit determinant.

If \(\lambda\) is a nonsquare modulo \(p\), the polynomial is
irreducible modulo \(p\). For any primitive vector \(v\),
\((v,Mv)\) is a basis modulo \(p\), hence over \(\mathcal O\),
and puts \(M\) in the companion form for \(t^2-\lambda\).
The same construction for \(L\) again gives \(g_0\).
Its centralizer algebra is
\(\mathcal O[t]/(t^2-\lambda)\), and determinant is the quadratic norm.
This norm is onto \(\mathcal O^\times\).
Indeed the residue norm is onto by cyclicity of the finite-field
multiplicative group. Lift a residue solution. The remaining unit
belongs to \(1+p\mathcal O\), and is the square of a scalar unit
by successive lifting, because two is invertible modulo \(p\).
The norm of that scalar unit is its square.
Thus also in this case the centralizer has every unit determinant.

Choose a centralizing \(c\) with
\(\det c=(\det g_0)^{-1}\) and put \(g=cg_0\).
Then \(g\in\operatorname{SL}_2(\mathcal O)\) and
\[
                   M=g^{-1}Lg,\qquad
                   z=g^{-1}\kappa(g).
\]
This proves (A.120.9) for every cocycle. The actual action in
(A.120.8) has \(\lambda=1\), so the split argument alone already
suffices for this application. The nonsplit argument records the
more general statement without an implicit square-root assumption.

\emph{A.120.3. The three-adic correction stays in the derived group.}
The actual three-adic coefficient is \(G_3'\), rather than \(G_3\).
The marked quotient of A.112 is
\[
       \chi:G_3\longrightarrow C_3,\qquad
       \ker\chi=G_3',\qquad \chi(X)=\chi(Y)=1.
 \tag{A.120.10}
\]
The original negative \(\alpha\) sends \(X\) to \(X^{-1}\)
and \(Y\) to a conjugate of \(Y^{-1}\), so it acts by \(-1\)
on \(\chi\). Conjugation by either word \(J_e^\alpha\)
acts trivially on this quotient: its pure factors act internally,
and \(s\) interchanges \(X,Y\), which have the same character.
Consequently the actual action, on all of \(G_3\), satisfies
\[
                         \chi(\kappa_e(g))=-\chi(g).
 \tag{A.120.11}
\]
This is not an assertion about an arbitrary matrix normalizer.

If \(z\in G_3'\) is a cocycle, (A.120.9) gives
\(g\in G_3\) with \(z=g^{-1}\kappa_e(g)\).
Applying (A.120.10)--(A.120.11) gives
\[
                         0=\chi(z)=-2\chi(g).
\]
Multiplication by two is invertible on \(C_3\), hence \(g\in G_3'\).
The same correcting element therefore belongs to the required
derived coefficient group, and
\[
                         H^1(C_2,G_3';\kappa_e)=\{*\}.
 \tag{A.120.12}
\]
There is no fixed-lifting assertion across a pro-two kernel.

\emph{A.120.4. Products, topology, and the entire forgetting map.}
The actual decomposition (A.120.1) is componentwise
\(\kappa_e\)-equivariant, as proved in A.112 and A.117.
Cocycles and correcting conjugators are consequently componentwise.
An arbitrary tuple of the conjugators just constructed is an element
of the compact product \(G_{\rm odd}'\).
Equations (A.120.9) and (A.120.12) prove that its entire cohomology
is trivial.

For a compact coefficient group with a continuous \(C_2\)-action,
the cocycle set is closed. Its equivalence relation is the orbit
relation of a compact group and is therefore closed.
The pointed quotient is compact Hausdorff.
The coordinate maps are continuous and induce the product
bijections on cohomology, hence homeomorphisms.
This argument requires no continuous choice of a cobounding
conjugator for each individual cocycle.

A.112 and A.117 compute the actual odd torus cohomology as the
product of its binary square-quotient factors in (A.120.4).
The inclusion \(T_{R,e}\hookrightarrow E_R'\) is exactly the
dyadic centralizer inclusion times these odd torus inclusions.
Their entire odd image is the basepoint by the preceding proof.
Thus the full map, with its topology, is exactly (A.120.5).
Replacing \(P'\) by \(P_*'\) gives the identical statement for
the earlier actual quotient \(E_1'\); all original odd coordinates
and frame words stay unchanged.

\emph{A.120.5. The constraint that the larger cohomology forgets.}
For an actual positive factor frame \(h\), the original defect is
\[
          w_h=(J_e^\alpha)^{-1}\alpha(h)^{-1}vh
                        \in T_{R,e}.
 \tag{A.120.13}
\]
Replacing \(h\) by \(ht\) preserves its positive-factor equation
exactly when \(t\in T_{R,e}\). Hence the coherence equation
must be solved in that centralizer, not in all of \(E_R'\).
The proof above shows that every odd class becomes an ambient
coboundary after this constraint is dropped. A correcting element
outside the centralizer changes the positive factor being framed;
it is not a coherent-frame correction.

In particular an odd centralizer class can be nontrivial while its
image in the larger coefficient group is the basepoint.
This theorem does not assert that the original full-GT class has
any specified value. It does not select \(\epsilon_A\), change
the original \(A_r\), or decide the original CM parity.
The full \(H_d\), the specified cubic norm gate of A.117,
and the transported pentagon of A.118 retain their stated scopes.
\end{proof}

\paragraph{Complete source record.}
The complete checkpoint35 proof is
\path{work_DE/odd_coherent_forgetting_obstruction35.md}.
Its entire independent reviews are
\path{work_A/review_odd_coherent_forgetting_obstruction35.md}
and
\path{work_CE/review_odd_coherent_forgetting_obstruction35.txt}.
The present proof prints the integral conjugacy, determinant correction,
three-adic derived correction, and all product-topology arguments.
It invokes no new external theorem, finite target enumeration,
numerical producer, or change of the original negative frame.
